% ===================== SLIDE 13 — FOUNDATIONS =====================
\begin{frame}[t]
\pmheader{FOUNDATIONS}{14 / 22}
\par\pmtitle{Information must be acquired, then aggregated}
\begin{pmbody}
{\color{pmGrey}\fontsize{20}{27}\selectfont Information must first be acquired and then incorporated into prices.\par}
\begin{itemize}\pmspreaditems
  \item Grossman and Stiglitz (1980): Fully revealing prices would eliminate the incentive to acquire costly information. Equilibrium must preserve compensation for becoming informed.
  \item Ostrovsky (2012): Within a dynamic model of strategic trading, separable securities aggregate dispersed information under specified conditions.
  \item Together, these studies connect information incentives with the conditions governing price formation.
\end{itemize}
\begin{pmcallout}[Potential contribution]
  Empirically connect changes in information costs with information acquisition, participation, and price informativeness.
\end{pmcallout}
\end{pmbody}
\end{frame}

% ===================== SLIDE 14 — TRADER COMPOSITION =====================
\begin{frame}[t]
\pmheader{TRADER COMPOSITION}{15 / 22}
\par\pmtitle{Trader composition may shape market accuracy}
\begin{pmbody}
{\color{pmGrey}\fontsize{20}{27}\selectfont Aggregate market performance may depend heavily on a small group of participants.\par}
\begin{itemize}\pmspreaditems
  \item Heesen and Luong (2026) identify skilled traders using earlier markets and examine their participation in separate test markets.
  \item A greater trading share of skilled switchers is associated with lower pricing error, primarily through better calibration.
  \item Akey et al. (2026) document concentrated profits on Polymarket and find that successful users more often provide liquidity.
\end{itemize}
\begin{pmcallout}[Interpretation]
  Who participates may matter alongside how much trading occurs.
\end{pmcallout}
\end{pmbody}
\end{frame}

% ===================== SLIDE 15 — MEASUREMENT =====================
\begin{frame}[t]
\pmheader{MEASUREMENT}{17 / 22}
\par\pmtitle{Activity, consensus, and accuracy are distinct}
\begin{pmbody}
{\color{pmGrey}\fontsize{20}{27}\selectfont Each outcome captures a different aspect of market performance.\par}
\begin{itemize}\pmspreaditems
  \item Trading volume measures activity.
  \item Liquidity measures describe trading conditions.
  \item Transaction-price dispersion captures differences in execution prices within a measurement window.
  \item Calibration evaluates whether realized event frequencies correspond to quoted probabilities.
\end{itemize}
\begin{pmcallout}[Page and Clemen (2013)]
  document horizon-dependent calibration and systematic favorite--longshot bias.
\end{pmcallout}
\begin{pmcallout}[Potential contribution]
  Explain when improvements in trading conditions or price consensus translate into better-calibrated forecasts.
\end{pmcallout}
\end{pmbody}
\end{frame}

% ===================== SLIDE 17 — CAUSAL GAP =====================
\begin{frame}[t]
\pmheader{CAUSAL GAP}{16 / 22}
\par\pmtitle{The central empirical gap is causal identification}
\begin{pmbody}
{\color{pmGrey}\fontsize{20}{27}\selectfont Existing studies describe trader behavior more clearly than they identify its effects.\par}
\begin{itemize}\pmspreaditems
  \item Skilled traders may choose markets that are already easier to price accurately.
  \item Identifying skill in earlier markets reduces circularity but does not eliminate market selection.
  \item Liquidity providers' profits may reflect execution advantages or compensation for risk, in addition to information advantages.
  \item These findings therefore do not establish that increasing skilled participation or maker activity causes better forecasts.
\end{itemize}
\begin{pmcallout}[Potential contribution]
  Identify a credible source of variation in participation and estimate its effect on market quality.
\end{pmcallout}
{\color{pmGrey}\fontsize{11}{14}\selectfont References: Akey et al. (2026); Heesen and Luong (2026).\par}
\end{pmbody}
\end{frame}
